% Copyright (C) 2012-2018 Dmitry V. Levin <ldv@altlinux.org>
% Permission is granted to copy, distribute and/or modify this document
% under the terms of the GNU Free Documentation License, Version 1.2
% or any later version published by the Free Software Foundation;
% with no Invariant Sections, no Front-Cover Texts, and no Back-Cover Texts.

\documentclass[unicode,aspectratio=169]{beamer}

\usepackage[utf8]{inputenc}
\usepackage[T2A]{fontenc}
\usepackage[english]{babel}
\usepackage{alltt}

\pgfdeclareimage[height=0.6cm]{basealt-corner-logo}{basealt.jpg}
\pgfdeclareimage[height=1.2cm]{strace-corner-logo}{strace-straus.pdf}
\pgfdeclareimage[height=3cm]{title-logo}{Logo_ikona_pruhledne.pdf}

\usetheme{Warsaw}
\setbeamertemplate{headline}{}
\setbeamertemplate{footline}{%
	\pgfuseimage{basealt-corner-logo}%
	\hfill%
	\pgfuseimage{strace-corner-logo}%
}
\setbeamertemplate{navigation symbols}{}
\setbeamertemplate{logo}{}
\setbeamercovered{transparent}

\title{\Huge Modern strace}
\author{\Huge Dmitry~Levin}
\date{\Large Brno, 2019}
\titlegraphic{\pgfuseimage{title-logo}}

\begin{document}

%%%%%%%
\begin{frame}
\titlepage
\end{frame}

%%%%%%%
\begin{frame}{Traditional strace \hfill [\insertframenumber/\inserttotalframenumber]}
\large
\begin{block}{Printing instruction pointer and timestamps}
\begin{itemize}
\item print instruction pointer: \textbf{-i} option
\item print timestamps: \textbf{-r}, \textbf{-t}, \textbf{-tt}, \textbf{-ttt}, and \textbf{-T} options
\end{itemize}
\end{block}

\begin{block}{Size and format of strings}
\begin{itemize}
\item string size: \textbf{-s} option
\item string format: \textbf{-x} and \textbf{-xx} options
\end{itemize}
\end{block}

\begin{block}{Verbosity of syscall decoding}
\begin{itemize}
\item abbreviate output: \textbf{-e abbrev=\textit{set}}, \textbf{-v} option
\item dereference structures: \textbf{-e verbose=\textit{set}}
\item print raw undecoded syscalls: \textbf{-e raw=\textit{set}}
\end{itemize}
\end{block}
\end{frame}

%%%%%%%
\begin{frame}{Traditional strace \hfill [\insertframenumber/\inserttotalframenumber]}
\large
\begin{block}{Printing signals}
\begin{itemize}
\item print signals: \textbf{-e signal=\textit{set}}
\end{itemize}
\end{block}

\begin{block}{Dumping}
\begin{itemize}
\item dump the data read from the specified descriptors: \textbf{-e read=\textit{set}}
\item dump the data written to the specified descriptors: \textbf{-e write=\textit{set}}
\end{itemize}
\end{block}

\begin{block}{Redirecting output to files or pipelines}
\begin{itemize}
\item write the trace to a file or pipeline: \textbf{-o \textit{filename}} option
\item write traces of processes to separate files: \textbf{-ff -o \textit{filename}}
\end{itemize}
\end{block}
\end{frame}

%%%%%%%
\begin{frame}{Traditional strace \hfill [\insertframenumber/\inserttotalframenumber]}
\large
\begin{block}{System call filtering}
\begin{itemize}
\item trace only the specified set of system calls: \textbf{-e trace=\textit{set}}
\end{itemize}
\end{block}

\begin{block}{System call statistics}
\begin{itemize}
\item count time, calls, and errors for each system call: \textbf{-c} option
\item sort the histogram printed by the \textbf{-c} option: \textbf{-S \textit{sortby}} option
\end{itemize}
\end{block}

\begin{block}{Tracing control}
\begin{itemize}
\item attach to existing processes: \textbf{-p \textit{pid}} option
\item trace child processes: \textbf{-f} option
\end{itemize}
\end{block}
\end{frame}

%%%%%%
\begin{frame}{Modern strace \hfill [\insertframenumber/\inserttotalframenumber]}
\large
\begin{block}{\large Tracing output format}
\begin{itemize}
\item pathnames accessed by name or descriptor: \textbf{-y} option
\item network protocol associated with descriptors: \textbf{-yy} option
\item stack of function calls: \textbf{-k} option
\item open \textbf{-o} in append mode: \textbf{-A} option
\item format of named constants and flags: \textbf{-X} option
\item elaborate syscall parsers
\end{itemize}
\end{block}

\begin{block}{\large System call filtering}
\begin{itemize}
\item pathnames accessed by name or descriptor: \textbf{-P} option
\item regular expressions: -e \textbf{trace}=/\textit{regexp}
\item optional specifications: -e \textbf{trace}=?\textit{spec}
\item new syscall classes: \%stat, \%lstat, \%fstat, \%statfs, \%fstatfs, \%\%stat, \%\%statfs
\end{itemize}
\end{block}
\end{frame}

%%%%%%%
\begin{frame}{Modern strace \hfill [\insertframenumber/\inserttotalframenumber]}
\large
\begin{block}{\large System call statistics}
\begin{itemize}
\item wall clock time spent in syscalls: \textbf{-w} option
\item combine statistics with regular output: \textbf{-C} option
\end{itemize}
\end{block}

\begin{block}{\large Tracing control}
\begin{itemize}
\item attach to multiple processes: \textbf{-p \textit{pid\_set}} option
\item detach on execve: \textbf{-b execve} option
\item run as a detached grandchild: \textbf{-D} option
\item interruptibility: \textbf{-I} option
\item postprocessing: \textbf{strace-log-merge}
\end{itemize}
\end{block}
\end{frame}

%%%%%%%
\begin{frame}{Modern strace \hfill [\insertframenumber/\inserttotalframenumber]}
\large
\begin{block}{\large System call tampering}
\begin{itemize}
\item fault injection: \\
-e \textbf{inject}=\textit{set}:\textbf{error}=\textit{errno}[:\textbf{when}=\textit{expr}][:\textbf{syscall}=\textit{syscall}]
\item return value injection: \\
-e \textbf{inject}=\textit{set}:\textbf{retval}=\textit{value}[:\textbf{when}=\textit{expr}][:\textbf{syscall}=\textit{syscall}]
\item signal injection: \\
-e \textbf{inject}=\textit{set}:\textbf{signal}=\textit{set}
\item delay injection: \\
-e \textbf{inject}=\textit{set}:\textbf{delay\_enter}=\textit{usecs} \\
-e \textbf{inject}=\textit{set}:\textbf{delay\_exit}=\textit{usecs}
\end{itemize}
\end{block}
\end{frame}

%%%%%%%
\begin{frame}[fragile]{System call filtering issues \hfill [\insertframenumber/\inserttotalframenumber]}
\begin{block}{\large glibc: open or openat?}
\begin{alltt}
glibc-2.25$ strace -qq \textcolor{red}{-e open} cat /dev/null
open("/etc/ld.so.cache", O_RDONLY|O_CLOEXEC) = 3
open("/lib64/libc.so.6", O_RDONLY|O_CLOEXEC) = 3
open("/dev/null", O_RDONLY)             = 3
glibc-2.26$ strace -qq \textcolor{red}{-e open} cat /dev/null
glibc-2.26$ strace -qq \textcolor{red}{-e openat} cat /dev/null
openat(AT_FDCWD, "/etc/ld.so.cache", O_RDONLY|O_CLOEXEC) = 3
openat(AT_FDCWD, "/lib64/libc.so.6", O_RDONLY|O_CLOEXEC) = 3
openat(AT_FDCWD, "/dev/null", O_RDONLY) = 3
glibc-2.25$ strace -qq \textcolor{red}{-e openat} cat /dev/null
\end{alltt}
\end{block}

\begin{block}{\large traditional approach is not portable}
\begin{alltt}
riscv$ strace \textcolor{red}{-e open,openat}
strace: invalid system call 'open'
\end{alltt}
\end{block}
\end{frame}

%%%%%%%
\begin{frame}[fragile]{regular expressions: -e trace=/\textit{regexp} \hfill [\insertframenumber/\inserttotalframenumber]}
\begin{block}{\large naive approach is inexact}
\large
\begin{alltt}
$ asinfo --set-arch x86_64,riscv --list-abi \textbackslash
  --nargs --get-sname \textcolor{red}{/open}
|   |                   | x86_64 | x86_64 | riscv | riscv |
| N |      Syscall name |  64bit |  32bit | 64bit | 32bit |
| 1 |           mq_open |      4 |      4 |     4 |     4 |
\textcolor{red}{| 2 |              open |      3 |      3 |     - |     - |}
| 3 | open_by_handle_at |      3 |      3 |     3 |     3 |
\textcolor{red}{| 4 |            openat |      4 |      4 |     4 |     4 |}
| 5 |   perf_event_open |      5 |      5 |     5 |     5 |
\end{alltt}
\textbf{asinfo} stands for \textbf{A}dvanced \textbf{S}yscall \textbf{INFO}rmation tool
\end{block}

\begin{block}{\large accurate and portable approach}
\texttt{\$ strace -e '/{\textasciicircum}open(at)?\$'}
\end{block}
\end{frame}

%%%%%%%
\begin{frame}[fragile]{New system call classes \hfill [\insertframenumber/\inserttotalframenumber]}
\begin{block}{\large traditional syscall classes}
\begin{itemize}
\item \textbf{desc}: take or return a descriptor
\item \textbf{file}: take a file name
\item \textbf{memory}: memory mapping, memory policy
\item \textbf{process}: process management
\item \textbf{signal}: signal related
\item \textbf{ipc}: SysV IPC related
\item \textbf{network}: network related
\end{itemize}
\vspace{-2\topsep}
\begin{alltt}
$ strace -e trace=\textit{class}
$ strace -e \textit{class}
\end{alltt}
\end{block}

\begin{block}{\large all syscall classes now have \% prefix}
\begin{alltt}
$ strace -e trace=%\textit{class}
$ strace -e %\textit{class}
\end{alltt}
\end{block}
\end{frame}

%%%%%%%
\begin{frame}[fragile]{New system call classes \hfill [\insertframenumber/\inserttotalframenumber]}
\begin{block}{\large new syscall classes}
\begin{itemize}
	\item \%stat, \%lstat, \%fstat
	\item \%statfs, \%fstatfs
	\item \%\%stat = \%stat + \%lstat + \%fstat + statx
	\item \%\%statfs = \%statfs + \%fstatfs + ustat
\end{itemize}
\end{block}

\begin{block}{\large strace \textbf{-y} -e \textbf{\%\%stat} ls /var/empty}
\begin{alltt}
\textcolor{red}{fstat}(3</etc/ld.so.cache>, {st_mode=S_IFREG|0644, st_size=30341, ...}) = 0
...
\textcolor{red}{fstat}(3</proc/filesystems>, {st_mode=S_IFREG|0444, st_size=0, ...}) = 0
\textcolor{red}{stat}("/var/empty", {st_mode=S_IFDIR|0555, st_size=40, ...}) = 0
\textcolor{red}{fstat}(3</var/empty>, {st_mode=S_IFDIR|0555, st_size=40, ...}) = 0
+++ exited with 0 +++
\end{alltt}
\end{block}
\end{frame}

%%%%%%%
\begin{frame}[fragile]{New system call classes \hfill [\insertframenumber/\inserttotalframenumber]}
\begin{block}{\large \%stat + \%lstat + \%fstat + statx = \%\%stat}
\begin{alltt}
$ asinfo --set-arch x86_64,riscv --list-abi --nargs --get-sname \%\%stat
|    |              | x86_64 | x86_64 | riscv | riscv |
|  N | Syscall name |  64bit |  32bit | 64bit | 32bit |
\textcolor{blue}{|  1 |        fstat |      2 |      2 |     2 |     - |}
\textcolor{blue}{|  2 |      fstat64 |      - |      2 |     - |     2 |}
\textcolor{blue}{|  3 |    fstatat64 |      - |      4 |     - |     4 |}
\textcolor{green}{|  4 |        lstat |      2 |      2 |     - |     - |}
\textcolor{green}{|  5 |      lstat64 |      - |      2 |     - |     - |}
\textcolor{blue}{|  6 |   newfstatat |      4 |      - |     4 |     - |}
\textcolor{blue}{|  7 |     oldfstat |      - |      2 |     - |     - |}
\textcolor{green}{|  8 |     oldlstat |      - |      2 |     - |     - |}
\textcolor{red}{|  9 |      oldstat |      - |      2 |     - |     - |}
\textcolor{red}{| 10 |         stat |      2 |      2 |     - |     - |}
\textcolor{red}{| 11 |       stat64 |      - |      2 |     - |     - |}
| 12 |        statx |      5 |      5 |     5 |     5 |
\end{alltt}
\end{block}
\end{frame}

%%%%%%%
\begin{frame}[fragile]{Network protocol associated with descriptors: \textbf{-yy} option \hfill [\insertframenumber/\inserttotalframenumber]}
\scriptsize
\begin{alltt}
$ strace \textcolor{red}{-yy} -e %desc,%network netcat 127.0.0.1 22 </dev/null
\ldots
socket(AF_INET, SOCK_STREAM, IPPROTO_TCP) = 3\textcolor{red}{<TCP:[518663]>}
connect(3\textcolor{red}{<TCP:[518663]>}, {sa_family=AF_INET, sin_port=htons(22),
        sin_addr=inet_addr("127.0.0.1")}, 16) = 0
poll([{fd=3\textcolor{red}{<TCP:[127.0.0.1:45678->127.0.0.1:22]>}, events=POLLIN},
     {fd=0\textcolor{red}{</dev/null<char 1:3>>}, events=POLLIN}], 2, -1) = 1 ([{fd=0, revents=POLLIN}])
read(0\textcolor{red}{</dev/null<char 1:3>>}, "", 2048)  = 0
shutdown(3\textcolor{red}{<TCP:[127.0.0.1:45678->127.0.0.1:22]>}, SHUT_WR) = 0
poll([{fd=3\textcolor{red}{<TCP:[127.0.0.1:45678->127.0.0.1:22]>}, events=POLLIN}, {fd=-1}], 2, -1)
     = 1 ([{fd=3, revents=POLLIN}])
read(3\textcolor{red}{<TCP:[127.0.0.1:45678->127.0.0.1:22]>}, "SSH-2.0-OpenSSH_7.9{\textbackslash}r{\textbackslash}n", 2048) = 21
write(1\textcolor{red}{</dev/pts/9<char 136:9>>}, "SSH-2.0-OpenSSH_7.9{\textbackslash}r{\textbackslash}n", 21) = 21
poll([{fd=3\textcolor{red}{<TCP:[127.0.0.1:45678->127.0.0.1:22]>}, events=POLLIN}, {fd=-1}], 2, -1)
     = 1 ([{fd=3, revents=POLLIN|POLLHUP}])
read(3\textcolor{red}{<TCP:[127.0.0.1:45678->127.0.0.1:22]>}, "", 2048) = 0
shutdown(3\textcolor{red}{<TCP:[127.0.0.1:45678->127.0.0.1:22]>}, SHUT_RD)
         = -1 ENOTCONN (Transport endpoint is not connected)
close(3\textcolor{red}{<TCP:[127.0.0.1:45678->127.0.0.1:22]>}) = 0
+++ exited with 0 +++
\end{alltt}
\end{frame}

%%%%%%%
\begin{frame}[fragile]{Stack of function calls: \textbf{-k} option \hfill [\insertframenumber/\inserttotalframenumber]}
\begin{block}{\large strace -qq -P /dev/full cat /dev/null > /dev/full}
\scriptsize
\begin{alltt}
fstat(1, {st_mode=S_IFCHR|0666, st_rdev=makedev(1, 7), ...}) = 0
close(1)                                = 0
\end{alltt}
\end{block}

\begin{block}{\large strace \textbf{-k} -qq -P /dev/full cat /dev/null > /dev/full}
\scriptsize
\begin{alltt}
fstat(1, {st_mode=S_IFCHR|0666, st_rdev=makedev(1, 7), ...}) = 0
 > /lib64/libc-2.27.so(__fxstat64+0x13) [0xe79c3]
 > /bin/cat(main+0x1b3) [0x4017e3]
 > /lib64/libc-2.27.so(__libc_start_main+0xe6) [0x21bd6]
 > /bin/cat(_start+0x29) [0x402179]
close(1)                                = 0
 > /lib64/libc-2.27.so(__close_nocancel+0x7) [0xe8b47]
 > /lib64/libc-2.27.so(_IO_file_close_it@@GLIBC_2.2.5+0x67) [0x79fd7]
 > /lib64/libc-2.27.so(fclose@@GLIBC_2.2.5+0x136) [0x6d376]
 > /bin/cat(close_stream+0x19) [0x404ce9]
 > /bin/cat(close_stdout+0x11) [0x402691]
 > /lib64/libc-2.27.so(__run_exit_handlers+0x170) [0x379c0]
 > /lib64/libc-2.27.so(exit+0x19) [0x37ab9]
 > /lib64/libc-2.27.so(__libc_start_main+0xed) [0x21bdd]
 > /bin/cat(_start+0x29) [0x402179]
\end{alltt}
\end{block}
\end{frame}

%%%%%%%
\begin{frame}[fragile]{Format of named constants and flags: \textbf{-X} option \hfill [\insertframenumber/\inserttotalframenumber]}
\begin{block}{\large strace -e /open cat /dev/null}
\scriptsize
\begin{alltt}
openat(\textcolor{red}{AT_FDCWD}, "/etc/ld.so.cache", \textcolor{red}{O_RDONLY|O_CLOEXEC}) = 3
openat(\textcolor{red}{AT_FDCWD}, "/lib64/libc.so.6", \textcolor{red}{O_RDONLY|O_CLOEXEC}) = 3
openat(\textcolor{red}{AT_FDCWD}, "/dev/null", \textcolor{red}{O_RDONLY}) = 3
+++ exited with 0 +++
\end{alltt}
\end{block}

\begin{block}{\large strace \textbf{-X verbose} -e /open cat /dev/null}
\scriptsize
\begin{alltt}
openat(\textcolor{red}{-100 /* AT_FDCWD */}, "/etc/ld.so.cache",
       \textcolor{red}{0x80000 /* O_RDONLY|O_CLOEXEC */}) = 3
openat(\textcolor{red}{-100 /* AT_FDCWD */}, "/lib64/libc.so.6",
       \textcolor{red}{0x80000 /* O_RDONLY|O_CLOEXEC */}) = 3
openat(\textcolor{red}{-100 /* AT_FDCWD */}, "/dev/null", \textcolor{red}{0 /* O_RDONLY */}) = 3
+++ exited with 0 +++
\end{alltt}
\end{block}

\begin{block}{\large strace \textbf{-X raw} -e /open cat /dev/null}
\scriptsize
\begin{alltt}
openat(\textcolor{red}{-100}, "/etc/ld.so.cache", \textcolor{red}{0x80000}) = 3
openat(\textcolor{red}{-100}, "/lib64/libc.so.6", \textcolor{red}{0x80000}) = 3
openat(\textcolor{red}{-100}, "/dev/null", \textcolor{red}{0})            = 3
+++ exited with 0 +++
\end{alltt}
\end{block}
\end{frame}

%%%%%%%
\begin{frame}[fragile]{System time spent in syscalls: \textbf{-c} option \hfill [\insertframenumber/\inserttotalframenumber]}
\begin{block}{\large strace \textbf{-c} sleep 1}
\small
\begin{alltt}
% time  seconds usecs/call calls errors syscall
------ -------- ---------- ----- ------ ----------
 31.45 0.000078         19     4        close
 25.40 0.000063         15     4        mprotect
\textcolor{red}{ 12.90 0.000032         32     1        nanosleep}
 11.29 0.000028          5     5        mmap
  8.47 0.000021          5     4        brk
  8.06 0.000020         20     1        munmap
  2.42 0.000006          6     1        arch_prctl
  0.00 0.000000          0     1        read
  0.00 0.000000          0     2        fstat
  0.00 0.000000          0     1      1 access
  0.00 0.000000          0     1        execve
  0.00 0.000000          0     2        openat
------ -------- ---------- ----- ------ ----------
100.00 0.000248               27      1 total
\end{alltt}
\end{block}
\end{frame}

%%%%%%%
\begin{frame}[fragile]{Wall clock time spent in syscalls: \textbf{-c -w} option \hfill [\insertframenumber/\inserttotalframenumber]}
\begin{block}{\large strace \textbf{-c -w} sleep 1}
\small
\begin{alltt}
% time  seconds usecs/call calls errors syscall
------ -------- ---------- ----- ------ ----------
\textcolor{red}{ 99.91 1.000184    1000183     1        nanosleep}
  0.04 0.000367        366     1        execve
  0.02 0.000216         53     4        close
  0.01 0.000087         17     5        mmap
  0.01 0.000075         18     4        mprotect
  0.01 0.000052         13     4        brk
  0.00 0.000037         18     2        openat
  0.00 0.000027         26     1        munmap
  0.00 0.000024         12     2        fstat
  0.00 0.000019         19     1      1 access
  0.00 0.000015         14     1        read
  0.00 0.000013         13     1        arch_prctl
------ -------- ---------- ----- ------ ----------
100.00 1.001116               27      1 total
\end{alltt}
\end{block}
\end{frame}

%%%%%%%
\begin{frame}[fragile]{Running as a detached grandchild: \textbf{-D} option \hfill [\insertframenumber/\inserttotalframenumber]}
\begin{block}{\large foreground strace}
\begin{alltt}
$ echo $$ && strace -e none sh -c 'echo $PPID'
1234
23456
+++ exited with 0 +++
$ echo $$ && strace -e none sh -c 'echo $PPID'
1234
23459
+++ exited with 0 +++
\end{alltt}
\end{block}

\begin{block}{\large background strace}
\begin{alltt}
$ echo $$ && strace \textcolor{red}{-D} -e none sh -c 'echo $PPID'
1234
1234
+++ exited with 0 +++
\end{alltt}
\end{block}
\end{frame}

%%%%%%%
\begin{frame}{Netlink socket parsers \hfill [\insertframenumber/\inserttotalframenumber]}
\tiny
\begin{block}{\large Currently supported netlink protocols}
\begin{itemize}
\item NETLINK\_AUDIT
\item NETLINK\_CRYPTO
\item NETLINK\_KOBJECT\_UEVENT
\item NETLINK\_NETFILTER
\item NETLINK\_ROUTE
\item NETLINK\_SELINUX
\item NETLINK\_SOCK\_DIAG
\item NETLINK\_XFRM
\item NETLINK\_GENERIC
\end{itemize}
\end{block}

\begin{block}{\large NETLINK\_ROUTE: ip route list table all}
\begin{alltt}
broadcast 127.0.0.0 dev lo table local proto kernel scope link src 127.0.0.1

local 127.0.0.0/8 dev lo table local proto kernel scope host src 127.0.0.1

local 127.0.0.1 dev lo table local proto kernel scope host src 127.0.0.1

broadcast 127.255.255.255 dev lo table local proto kernel scope link src 127.0.0.1
\end{alltt}
\end{block}
\end{frame}

%%%%%%%
\begin{frame}{Netlink socket parsers: NETLINK\_ROUTE \hfill [\insertframenumber/\inserttotalframenumber]}
\tiny
\begin{block}{\large strace -e trace=sendto,recvmsg ip route list}
\begin{alltt}
sendto(3, \{\{len=40, type=RTM\_GETROUTE, flags=NLM\_F\_REQUEST|NLM\_F\_DUMP, seq=1357924680, pid=0\},
 \textcolor{blue}{\{rtm\_family=AF\_UNSPEC, rtm\_dst\_len=0, rtm\_src\_len=0, rtm\_tos=0, rtm\_table=RT\_TABLE\_UNSPEC, rtm\_protocol=RTPROT\_UNSPEC, rtm\_scope=RT\_SCOPE\_UNIVERSE, rtm\_type=RTN\_UNSPEC, rtm\_flags=0\}},
  \textcolor{red}{\{nla\_len=0, nla\_type=RTA\_UNSPEC\}}\}, 40, 0, NULL, 0) = 40

\smallskip
recvmsg(3, \{msg\_name=\{sa\_family=AF\_NETLINK, nl\_pid=0, nl\_groups=00000000\}, msg\_namelen=12, msg\_iov=[\{iov\_base=[
\textcolor{brown}{\{\{len=60, type=RTM\_NEWROUTE, flags=NLM\_F\_MULTI, seq=1357924680, pid=12345\},
 \textcolor{blue}{\{rtm\_family=AF\_INET, rtm\_dst\_len=32, rtm\_src\_len=0, rtm\_tos=0, rtm\_table=RT\_TABLE\_LOCAL, rtm\_protocol=RTPROT\_KERNEL, rtm\_scope=RT\_SCOPE\_LINK, rtm\_type=RTN\_BROADCAST, rtm\_flags=0\}},
  \textcolor{red}{[\{\{nla\_len=8, nla\_type=RTA\_TABLE\}, RT\_TABLE\_LOCAL\}, \{\{nla\_len=8, nla\_type=RTA\_DST\}, inet\_addr("127.0.0.0")\}, \{\{nla\_len=8, nla\_type=RTA\_PREFSRC\}, inet\_addr("127.0.0.1")\}, \{\{nla\_len=8, nla\_type=RTA\_OIF\}, if\_nametoindex("lo")\}]}\}},
\textcolor{brown}{\{\{len=60, type=RTM\_NEWROUTE, flags=NLM\_F\_MULTI, seq=1357924680, pid=12345\},
 \textcolor{blue}{\{rtm\_family=AF\_INET, rtm\_dst\_len=8, rtm\_src\_len=0, rtm\_tos=0, rtm\_table=RT\_TABLE\_LOCAL, rtm\_protocol=RTPROT\_KERNEL, rtm\_scope=RT\_SCOPE\_HOST, rtm\_type=RTN\_LOCAL, rtm\_flags=0\}},
  \textcolor{red}{[\{\{nla\_len=8, nla\_type=RTA\_TABLE\}, RT\_TABLE\_LOCAL\}, \{\{nla\_len=8, nla\_type=RTA\_DST\}, inet\_addr("127.0.0.0")\}, \{\{nla\_len=8, nla\_type=RTA\_PREFSRC\}, inet\_addr("127.0.0.1")\}, \{\{nla\_len=8, nla\_type=RTA\_OIF\}, if\_nametoindex("lo")\}]}\}},
\textcolor{brown}{\{\{len=60, type=RTM\_NEWROUTE, flags=NLM\_F\_MULTI, seq=1357924680, pid=12345\},
 \textcolor{blue}{\{rtm\_family=AF\_INET, rtm\_dst\_len=32, rtm\_src\_len=0, rtm\_tos=0, rtm\_table=RT\_TABLE\_LOCAL, rtm\_protocol=RTPROT\_KERNEL, rtm\_scope=RT\_SCOPE\_HOST, rtm\_type=RTN\_LOCAL, rtm\_flags=0\}},
  \textcolor{red}{[\{\{nla\_len=8, nla\_type=RTA\_TABLE\}, RT\_TABLE\_LOCAL\}, \{\{nla\_len=8, nla\_type=RTA\_DST\}, inet\_addr("127.0.0.1")\}, \{\{nla\_len=8, nla\_type=RTA\_PREFSRC\}, inet\_addr("127.0.0.1")\}, \{\{nla\_len=8, nla\_type=RTA\_OIF\}, if\_nametoindex("lo")\}]}\}},
\textcolor{brown}{\{\{len=60, type=RTM\_NEWROUTE, flags=NLM\_F\_MULTI, seq=1357924680, pid=12345\},
 \textcolor{blue}{\{rtm\_family=AF\_INET, rtm\_dst\_len=32, rtm\_src\_len=0, rtm\_tos=0, rtm\_table=RT\_TABLE\_LOCAL, rtm\_protocol=RTPROT\_KERNEL, rtm\_scope=RT\_SCOPE\_LINK, rtm\_type=RTN\_BROADCAST, rtm\_flags=0\}},
  \textcolor{red}{[\{\{nla\_len=8, nla\_type=RTA\_TABLE\}, RT\_TABLE\_LOCAL\}, \{\{nla\_len=8, nla\_type=RTA\_DST\}, inet\_addr("127.255.255.255")\}, \{\{nla\_len=8, nla\_type=RTA\_PREFSRC\}, inet\_addr("127.0.0.1")\}, \{\{nla\_len=8, nla\_type=RTA\_OIF\}, if\_nametoindex("lo")\}]}\}}
], iov\_len=32768\}], msg\_iovlen=1, msg\_controllen=0, msg\_flags=0\}, 0) = 240 \\
\ldots
\end{alltt}
\end{block}
\end{frame}

%%%%%%%
\begin{frame}[fragile]{Postprocessing: \textbf{strace-log-merge} \hfill [\insertframenumber/\inserttotalframenumber]}
\scriptsize
\begin{alltt}
$ strace -o log -ff -tt -e trace=execve,nanosleep sh -c 'sleep 0.1 & sleep 0.2 & sleep 0.3'
$ strace-log-merge log
13475 21:13:52.040837 execve("/bin/sh", ["sh", "-c", "sleep 0.1 & sleep 0.2 & sleep 0."...],
      0x7ffde54b2450 /* 33 vars */) = 0
\textcolor{blue}{13478 21:13:52.044050 execve("/bin/sleep", ["sleep", "0.3"], 0x5631be4f87a8 /* 33 vars */) = 0}
\textcolor{red}{13476 21:13:52.044269 execve("/bin/sleep", ["sleep", "0.1"], 0x5631be4f87a8 /* 33 vars */) = 0}
\textcolor{green}{13477 21:13:52.044389 execve("/bin/sleep", ["sleep", "0.2"], 0x5631be4f87a8 /* 33 vars */) = 0}
\textcolor{blue}{13478 21:13:52.046207 nanosleep(\{tv_sec=0, tv_nsec=300000000\}, NULL) = 0}
\textcolor{red}{13476 21:13:52.046303 nanosleep(\{tv_sec=0, tv_nsec=100000000\}, NULL) = 0}
\textcolor{green}{13477 21:13:52.046318 nanosleep(\{tv_sec=0, tv_nsec=200000000\}, NULL) = 0}
\textcolor{red}{13476 21:13:52.146852 +++ exited with 0 +++}
13475 21:13:52.146942 --- SIGCHLD \{si_signo=SIGCHLD, si_code=CLD_EXITED,
      si_pid=13476, si_uid=1000, si_status=0, si_utime=0, si_stime=0\} ---
\textcolor{green}{13477 21:13:52.247782 +++ exited with 0 +++}
13475 21:13:52.247885 --- SIGCHLD \{si_signo=SIGCHLD, si_code=CLD_EXITED,
      si_pid=13477, si_uid=1000, si_status=0, si_utime=0, si_stime=0\} ---
\textcolor{blue}{13478 21:13:52.347680 +++ exited with 0 +++}
13475 21:13:52.347786 --- SIGCHLD \{si_signo=SIGCHLD, si_code=CLD_EXITED,
      si_pid=13478, si_uid=1000, si_status=0, si_utime=0, si_stime=0\} ---
13475 21:13:52.348069 +++ exited with 0 +++
\end{alltt}
\end{frame}

%%%%%%%
\begin{frame}[fragile]{System call tampering: fault injection \hfill [\insertframenumber/\inserttotalframenumber]}
\begin{block}{\large -e \textbf{inject}=\textit{set}:\textbf{error}=\textit{errno}[:\textbf{when}=\textit{expr}][:\textbf{syscall}=\textit{syscall}]}
\textbf{inject}=\textit{set} -- fault injection for the specified set of syscalls \\
\textbf{error}=\textit{errno} -- the error code to fail syscalls with \\
\textbf{when}=\textit{expr} -- when to inject, in \textit{first[+[step]]} form \\
\textbf{syscall}=\textit{syscall} -- inject the specified syscall instead of -1
\end{block}

\begin{block}{\large strace -e /open -e \textbf{inject}=all:\textbf{error}=EACCES:\textbf{when}=3 \textbackslash\\
cat /dev/full /dev/null}
\begin{alltt}
openat(AT_FDCWD, "/etc/ld.so.cache", O_RDONLY|O_CLOEXEC) = 3
openat(AT_FDCWD, "/lib64/libc.so.6", O_RDONLY|O_CLOEXEC) = 3
\textcolor{red}{openat(AT_FDCWD, "/dev/full", O_RDONLY)
 = -1 EACCES (Permission denied) (INJECTED)}
cat: /dev/full: Permission denied
openat(AT_FDCWD, "/dev/null", O_RDONLY) = 3
+++ exited with 1 +++
\end{alltt}
\end{block}
\end{frame}

%%%%%%%
\begin{frame}[fragile]{System call tampering: fault injection \hfill [\insertframenumber/\inserttotalframenumber]}
\begin{block}{\large python3.5 bug: error opening /dev/urandom}
\small
\begin{alltt}
$ strace -P /dev/urandom -e inject=\%file:error=ENOENT python3
\textcolor{red}{openat(AT_FDCWD, "/dev/urandom", O_RDONLY|O_CLOEXEC)
 = -1 ENOENT (No such file or directory) (INJECTED)}
Fatal Python error: Failed to open /dev/urandom
--- SIGSEGV \{si_signo=SIGSEGV, si_code=SEGV_MAPERR, si_addr=0x50\} ---
+++ killed by SIGSEGV +++
Segmentation fault
\end{alltt}
\end{block}

\begin{block}{\large python3.5 bug: error reading /dev/urandom}
\small
\begin{alltt}
$ strace -a0 -e read -P /dev/urandom -e inject=all:error=EIO python3
\textcolor{red}{read(3, 0x8db610, 24) = -1 EIO (Input/output error) (INJECTED)}
Fatal Python error: Failed to read bytes from /dev/urandom
--- SIGSEGV \{si_signo=SIGSEGV, si_code=SEGV_MAPERR, si_addr=0x50\} ---
+++ killed by SIGSEGV +++
Segmentation fault
\end{alltt}
\end{block}
\end{frame}

%%%%%%%
\begin{frame}[fragile]{System call tampering: fault injection \hfill [\insertframenumber/\inserttotalframenumber]}
\begin{block}{\large glibc <= 2.25 dynamic linker bug}
\small
\begin{alltt}
$ strace -e mprotect -efault=all:error=EPERM:when=1 pwd
\textcolor{red}{mprotect(0x7fabcd00f000, 2097152, PROT_NONE)
 = -1 EPERM (Operation not permitted) (INJECTED)}
mprotect(0x7fabcd20f000, 16384, PROT_READ) = 0
mprotect(0x606000, 4096, PROT_READ)     = 0
mprotect(0x7fabcd441000, 4096, PROT_READ) = 0
/
+++ exited with 0 +++
\end{alltt}
\end{block}

\begin{block}{\large glibc >= 2.26: with a proper check}
\small
\begin{alltt}
$ strace -e mprotect -efault=all:error=EPERM:when=1 pwd
\textcolor{red}{mprotect(0x7fabcd00f000, 2097152, PROT_NONE)
 = -1 EPERM (Operation not permitted) (INJECTED)}
pwd: error while loading shared libraries: libc.so.6:
 cannot change memory protections
+++ exited with 127 +++
\end{alltt}
\end{block}
\end{frame}

%%%%%%%
\begin{frame}[fragile]{System call tampering: return value injection \hfill [\insertframenumber/\inserttotalframenumber]}
\begin{block}{\large -e \textbf{inject}=\textit{set}:\textbf{retval}=\textit{value}[:\textbf{when}=\textit{expr}][:\textbf{syscall}=\textit{syscall}]}
\textbf{inject}=\textit{set} -- fault injection for the specified set of syscalls \\
\textbf{retval}=\textit{value} -- the return value to return \\
\textbf{when}=\textit{expr} -- when to inject, in \textit{first[+[step]]} form \\
\textbf{syscall}=\textit{syscall} -- inject the specified syscall instead of -1
\end{block}

\begin{block}{\large example: recovery of temporary files}
\begin{alltt}
$ cat script.sh
t=`mktemp`; trap 'rm -f "$t"' 0; echo secret $$ > $t
$ strace -qq -f -e signal=none -e /unlink
  -e inject=all:retval=0 sh script.sh
[pid 347] \textcolor{red}{unlinkat(AT_FDCWD, "/tmp/tmp.l1AlwyCYH3", 0) = 0 (INJECTED)}
$ cat /tmp/tmp.l1AlwyCYH3
secret 345
\end{alltt}
\end{block}

\end{frame}

%%%%%%%
\begin{frame}[fragile]{System call tampering: delay injection \hfill [\insertframenumber/\inserttotalframenumber]}
\begin{block}{\large syscall delay injection}
strace -e \textbf{inject}=\textit{set}:\textbf{delay\_enter}=\textit{usecs} \\
strace -e \textbf{inject}=\textit{set}:\textbf{delay\_exit}=\textit{usecs}
\end{block}

\begin{block}{\large dd if=/dev/zero of=/dev/null bs=1M count=10}
\begin{alltt}
10+0 records in
10+0 records out
10485760 bytes (10 MB, 10 MiB) copied, \textcolor{red}{0.00211354} s, \textcolor{red}{5.0 GB}/s
\end{alltt}
\end{block}

\begin{block}{\large strace -einject=write:delay\_exit=100000 -ewrite -o/dev/null \textbackslash\\
dd if=/dev/zero of=/dev/null bs=1M count=10}
\begin{alltt}
10+0 records in
10+0 records out
10485760 bytes (10 MB, 10 MiB) copied, \textcolor{red}{1.10658} s, \textcolor{red}{9.5 MB}/s
\end{alltt}
\end{block}
\end{frame}

%%%%%%%
{
\setbeamertemplate{footline}{}
\begin{frame}{Questions?}
	\begin{columns}
		\column{7cm}
\begin{block}{\large homepage}
	https://strace.io
\end{block}
\begin{block}{\large strace.git}
	https://github.com/strace/strace.git

	https://gitlab.com/strace/strace.git
\end{block}
\begin{block}{\large mailing list}
	strace-devel@lists.strace.io
\end{block}
\begin{block}{\large IRC channel}
	\#strace@freenode
\end{block}
		\column{3cm}
			\centerline{\includegraphics[height=7cm]{strace-straus.pdf}}
	\end{columns}
\end{frame}
}

\end{document}
